\section{单圈修正的结构}

\subsection{禁闭与红外截断}

计算 QCD 中胶子辐射修正有若干标准技术。它们大多面向绝对微扰 QCD（pQCD）适用区域。然而，在涉及不太大能标的应用中，直接使用这些方法可能需要小心。为此，让我们讨论一些处于微扰方法适用边界上的计算的特点。

首先应当记住，夸克与胶子是禁闭的，即所有图形的传播子（即使在连续情形）都嵌入在一个体积有限的区域中，其尺度由强子半径决定。禁闭效应特别地导致 ITD 或伪 PDF 之类的关联函数在超过强子半径 $R$ 的距离 $z_3$ 处快速下降。尽管如此，在短距离处仍可使用渐近自由，并由此得到 $\ln z_3^2$ 奇异性。

因此，把伪 ITD 与伪 PDF 视为软部与硬部之和是有意义的。软部基本反映系统的尺寸，并假定在 $z_3=0$ 时有限。硬部在 \mbox{$z_3 \to 0$} 时奇异，由微扰相互作用产生。于是可以把硬部可视化为由软部通过硬交换核 $H (0,z; z_1, z_2)$ 生成：
\begin{align}
    & {\cal M}^{\rm   hard} (\nu, -z^2=z_3^2)
    \nn & =   \int d^4z_1 \, d^4 z_2 \, H (0,z; z_1, z_2) \,  {\cal M}^{\rm  soft} (z_1,z_2)\  .
\label{Mhar}
\end{align}

在标准 pQCD 因子化方法中，软部用壳上部分子态模拟，而 \mbox{$\ln z_3^2$ 奇异性}或者表现为 $\ln (z_3^2 m^2)$（$m$ 为部分子质量），或者表现为 $\ln (z_3^2 \mu_{\rm IR}^2)$（$\mu_{\rm IR}$ 是无质量部分子情形下用于红外奇异性维数正规化的标度）。

由于式 (\ref{Mhar}) 中的 ${\cal  M}^{\rm  soft} (z_1,z_2)$ 在大间隔 $|z_1-z_2|$ 下迅速下降，即使夸克无质量，强子尺度 $R$ 也为积分提供了红外截断。虽然在短距离处得到 $\ln (z_3^2 / R^2)$ 行为，但对数形式只是对 $z_3 \ll R$ 成立的近似。这样的限制在格点上可能难以满足。

当然，禁闭所施加的红外正规化的确切形式是未知的。为获得直观感受，让我们取一个以质量项实现的红外正规化。产生 $\ln z_3^2$ 奇异性的一个典型积分形如
   \begin{align}
  I_K (z_3^2) = \int_0^\infty \frac{d\alpha}{\alpha} e^{-z_3^2/4\alpha -   \alpha m^2} \ ,
  \label{IK}
     \end{align}
其中 $\alpha$ 是 Schwinger 的 $\alpha$ 参数，$m$ 是红外调节量。可以看到
   \begin{align}
  I_K (z_3^2) = &2 K_0 (mz_3) \nn & = - \ln (m^2 z_3^2) + 2 \ln (2e^{-  \gamma_E })+ {\cal O} (z_3^2)  \  ,
  \label{Kexp}
     \end{align}
其中 $K_0(mz_3)$ 是修正贝塞尔函数。它在小 $z_3$ 处的展开明确显示出预期的 $\ln (z_3^2 m^2)$ 奇异性。

通常的 pQCD 因子化做法是把 $\ln (z_3/ R)$ 拆成归于系数函数的短距部分 $\ln (z_3 \mu)$ 与被吸收进``重正化''PDF $f(x,\mu^2)$ 的长距部分 $\ln ( 1/ \mu R )$。鉴于常用的格距 $a\sim 0.1$ fm 和强子尺度 $R \lesssim 1$ fm，问题在于 \mbox{$z_3$ 依赖}的对数部分是否还有足够的区间能在数据中显现出来。

贝塞尔函数 $K_0 (mz_3)$ 的一个重要特点是：当 $z_3$ 超过红外截断 $1/m$ 时它指数式衰减。因此，如果不用 $\ln (1/z_3^2)$ 这一 $I_K(z_3^2)$ 的短距近似，而将``精确''函数 $I_K(z_3^2)$ 用作演化项，那么大 $z_3$ 处将不存在演化修正。换言之，对数演化在大距离处消失。

\subsection{重标度关系}

为确定伪 PDF 标度 $z_3$ 与 $\overline{\rm MS}$ 标度 $\mu$ 之间的关系，必须计入常数项，如式 (\ref{Kexp}) 中的 $2 \ln (2e^{-  \gamma_E })$。在 $\overline{\rm MS}$-OPE 方法中，取 $z^2=0$，然后应用维数正规化——它给积分 (\ref{IK}) 加上 $\alpha^{\epsilon}$ 因子使之收敛。随后使用 $\overline{\rm MS}$ 约定，该约定在此情形下安排得恰好给出 $\ln (\mu^2/m^2)$ 作为结果：
   \begin{align}
  I_{Dm} (\mu^2) =& \int_0^\infty \frac{d\alpha}{\alpha}  (\alpha \mu^2 e^{\gamma_E})^{\epsilon}  e^{-   \alpha m^2} \,
  =\Gamma (\epsilon) \left ( \frac{\mu^2 e^{\gamma_E}}{m^2} \right )^\epsilon \nn & \to
  \frac1{\epsilon}   + \ln (\mu^2/m^2)  \  .
  \label{ImM}
     \end{align}
因此，式 (\ref{Kexp}) 中的常数项给出了伪 PDF 与 $\overline{\rm MS}$ 部分子分布之间由式 (\ref{Ptof}) 表达的领头对数重标度系数 $2e^{-  \gamma_E}$。

有人或许会问：若采用另一类红外正规化会怎样？特别是，TMD 的高斯模型暗示大 $z_3$ 处的下降也是高斯型的。可以预期硬修正在大 $z_3$ 处应类似于软部的行为。这样看来，修正贝塞尔函数的指数 $e^{-m|z_3|}$ 衰减可能显得太慢。对式 (\ref{IK}) 施加尖锐红外截断
    \begin{align}
  I_G (z_3^2) = \int_0^{z_0^2/4}  \frac{d\alpha}{\alpha}    e^{-z_3^2/4\alpha } = \Gamma [0,z_3^2/z_0^2]
  \label{IDexp}
     \end{align}
即可轻易给出高斯型下降。对小 $z_3^2$，不完全伽马函数 $\Gamma (0,z_3^2/z_3^2)$ 具有对数奇异性
   \begin{align}
  \Gamma (0,z_3^2/z_0^2) =  \ln (z_0^2/z_3^2) -  \gamma_E + {\cal O} (z_3^2)  \  ,
  \label{Dexp0}
     \end{align}
     而对大 $z_3^2$，函数 $I_G (z_3^2)$ 具有高斯型 $e^{-z_3^2/z_0^2}$ 衰减。我们同样可以用 $\overline{\rm MS}$ 方案计算式 (\ref{IDexp}) 的 $z_3=0$ 版本：
  \begin{align}
  I_{DG} (\mu^2) = & \int_0^{z_0^2/4}  \frac{d\alpha}{\alpha}  (\alpha \mu^2  e^{\gamma_E})^{\epsilon}    =
  \frac1{\epsilon} \,  \left (\frac{ z_0^2  \mu^2  e^{\gamma_E}}{4}  \right )^{\epsilon} \nn &
  \to  \frac1{\epsilon} \, + \ln (z_0^2 \mu^2    ) -2  \ln ( 2 e^{-\gamma_E}) -  \gamma_E \ .
  \label{IDexp2}
     \end{align}
可以看到伪 PDF/PDF 重标度关系 (\ref{Ptof}) 保持不变。这是自然的结果，因为有限 $z_3$ 截断与 $\overline{\rm MS}$ 截断之间的关系只涉及双定域算符的短距性质。

\subsection{单圈修正}

上一节的讨论只涉及两种正规化方案之间的整体重标度（正如 MOM 方案与 $\overline{\rm MS}$ 方案中 QCD 标度 $\Lambda$ 取值之间的关系）。要建立伪 PDF 与 $\overline{\rm MS}$-PDF 的联系，还需要式 \mbox{(\ref{OPE})} 中非局域 OPE 单圈系数函数的常数部分。文献 \cite{Ji:2017rah} 与 \cite{Radyushkin:2017lvu} 给出了结果，二者间略有差异。在复核计算并改正笔误之后，我们把结果写成
  \begin{align}
  &
{\mathfrak M}  (\nu, z_3^2)    = {\mathfrak  M}^{\rm soft}  ( \nu, 0)  -  \frac{\alpha_s}{2\pi} \, C_F
\int_0^1  dw \,  \left \{   \frac{1+w^2} {1-w}    \right.  \nn & \left.  \hspace{2cm} \times
  \left [  \ln \left (z_3^2m^2\frac{  e^{2\gamma_E}}{4} \right )   +1 \right ]
\right. \nn    & \left.
+ 4  \,  \frac{\ln (1-w)}{1-w}
\right \}  \left [
{\mathfrak  M}^{\rm soft}  (w \nu,0) -  {\mathfrak  M}^{\rm soft}  ( \nu,0) \right ]\  .
 \label{Mhard}
 \end{align}

转向 PDF 一侧，取 $z^2=0$ 并用 $\overline{\rm MS}$ 方案处理 UV 发散，得到
   \begin{align}
{\cal I}  (\nu, \mu^2)
  =&  {\mathfrak  M}^{\rm soft}  ( \nu, 0) \nn & -  \frac{\alpha_s}{2\pi} \, C_F\
    \int_0^1  dw \
  \left [
{\mathfrak  M}^{\rm soft}  (w \nu, 0) -  {\mathfrak  M}^{\rm soft}  ( \nu, 0) \right ]\nn &
  \nn & \times \left \{    \frac{1+w^2} {1-w} \,  \ln (m^2/ \mu^2)   + 2 (1-w) \right \}
.
\label{Msbar}
 \end{align}

此处的对数部分包含一个卷积，可符号化地写作 $B \otimes {\mathfrak M}\,  (\nu)$，其中
     \begin{align}
B(w) =
\left [ \frac{1+w^2}{1- w} \, \right ]_+
\label{AP}
 \end{align}
是 Altarelli-Parisi 劈裂核 \cite{Altarelli:1977zs}。

合并式 (\ref {Mhard}) 与 (\ref{Msbar})，得到关系
  \begin{align}
{\cal I}  (\nu, \mu^2)    = & {\mathfrak  M} ( \nu, z_3^2) +  \frac{\alpha_s}{2\pi} \, C_F\
 \int_0^1  dw \,   {\mathfrak  M}  (w \nu,z_3^2)   \nn &
 \times  \left  \{   B(w) \, \left [ \ln \left (z_3^2\mu^2 \frac{  e^{2\gamma_E}}{4} \right )  +1
 \right ]   \right. \nn & \left.
+   \left [ 4  \frac{\ln (1-w)}{1-w} - 2 (1-w)  \right ]_+ \right  \}
\
\label{MNL}
 \end{align}
它与文献 \cite{Izubuchi:2018srq} 的近期结果一致（另见文献 \cite{Zhang:2018ggy}）。式 (\ref{MNL}) 使人们能把 ${\cal  M} ( \nu, z_3^2) $ 的数据点转换为 ${\cal  I} (\nu, \mu^2)$ 的``数据''。

第二行的第一项显然反映 $z^2$ 截断与 $\overline{\rm MS}$ 截断之间一般的乘性标度差。若忽略其余各项，则 $ {\mathfrak  M} ( \nu, z_3^2) $ 与 ${\cal  I} (\nu, \mu^2) $ 之间的差别就只是重标度 \mbox{$\mu^2  =4e^{-2\gamma_E}/ z_3^2$}。此时可以把 $ {\mathfrak  M} ( \nu, z_3^2) $ 数据演化到某个特定值 $z_0$，并（在此近似下）把所得函数 $ {\mathfrak  M} ( \nu, z_0^2) $ 当作对应标度 \mbox{$\mu = 2 e^{-\gamma_E}/ z_0$}（数值上接近 $1/z_0$）的 $\overline{\rm MS}$ ITD。

这一简单重标度关系（用于文献 \cite{Orginos:2017kos}）在纳入式 (\ref{MNL}) 的其余项后被修改。特别是，正比于 Altarelli-Parisi 劈裂核 $B(w)$ 的项可被吸收进 $\ln z_3^2$ 项中，这只会把重标度关系改写为 $\mu = 2 e^{-1/2-\gamma_E}/ z_0$。

\begin{figure}[tbp]
    \centerline{\includegraphics[width=3in]{images/linkboth}}
    \caption{产生大单圈修正的图形的坐标表示。
        \label{link}}
    \end{figure}

含 $[\ln (1-w)]/(1-w)$ 的项产生一个大的负贡献。按文献 \cite{Radyushkin:2017lvu}，在费曼规范下它来自涉及规范连接的顶点图的演化部分（见图~\ref{link}）。关键在于胶子连接到连接上的一个跑动位置 $tz_3$。对 $t$ 积分等运算后，净结果是这些图形的 $z_3$ 依赖由小于 $z_3$ 的有效标度生成。事实上，让我们把 $[\ln (1-w)]/(1-w)$ 项与 $\ln z_3^2$ 对数合并，把式 (\ref{MNL}) 改写为
    \begin{align}
{\cal I}  (\nu, \mu^2)    = & {\mathfrak  M} ( \nu, z_3^2) +  \frac{\alpha_s}{\pi} \, C_F\
 \int_0^1  dw \,   {\mathfrak  M}  (w \nu,z_3^2)   \nn &
 \times  \Biggl   \{   \frac{1+w^2}{1-w}  \,  \ln \left [(1-w)z_3  \mu \frac{  e^{\gamma_E+1/2}}{2} \right ]
\nn &
+   \left [   (w+1)  \ln (1-w)  -(1-w) \right ] \Biggr  \} _+
\   .
\label{MNL2}
 \end{align}
我们看到 $z_3$ 现在通过跑动位置 $(1-w)z_3$ 进入。剩余的 $(w+1) \ln (1-w)$ 项在 $w=1$ 处远不如 $B(w)$ 奇异，不产生大贡献。

因此，单圈修正的大小由合并后的演化对数支配。由于它依赖于积分变量 $w$，无法通过选取特定的 $\mu$ 使之为零。不过，$w$ 积分后的贡献将在某个 $\mu$ 处消失，我们可把它写成
     \begin{align}
\mu =  \frac{2 e^{-1/2-\gamma_E}}{\langle 1-w \rangle   }\, \frac1{z_3} \ ,
\end{align}
其中 $\langle 1-w \rangle $ 是 $1-w$ 的``平均''值。既然 $B(w)$ 在 $w=1$ 处强烈增强，应预期 $\langle 1-w \rangle $ 数值很小，从而给出带相当大系数 $k$ 的 $\mu \sim k/z_3$ 型重标度。我们将看到，此处 $k \sim 4$。

再次可以问：含 $\ln z_3^2 $ 对数的微扰公式 (\ref{MNL}) 是否适用于实际的格点数据？特别是，我们关于质量项红外正规化及相应贝塞尔函数的练习表明，硬项的对数行为 $\ln z_3^2 $ 只在 $z_3 $ 远低于红外截断 $R$（在我们的情形即强子尺度）时才成立。因此，实际问题是：数据是否确实在某个\mbox{小 $z_3$}区域内显示出对数演化行为？
